\begin{theorem}[Adjunction unit-counit carrier-swap involution]
\label{thm:adjunction-unit-counit-carrier-swap-involution}
Let $D_i=(p_i,q_i,a_i,u_i,c_i,l_i,r_i)$ for $i=0,1,2$ be unit-counit displays
of histories. If
$$
\begin{array}{l}
\mathsf{AdjunctionUnitCounitCarrier}(p_0,q_0,a_0,u_0,c_0,l_0,r_0)\\
{}\land\mathsf{AdjunctionUnitCounitCarrierSwap}(D_0,D_1)
\land\mathsf{AdjunctionUnitCounitCarrierSwap}(D_1,D_2),
\end{array}
$$
then
$$
\begin{array}{l}
\mathsf{AdjunctionUnitCounitCarrier}(p_2,q_2,a_2,u_2,c_2,l_2,r_2)\\
{}\land\mathsf{AdjunctionUnitCounitDisplaySame}(D_0,D_2).
\end{array}
$$
\end{theorem}
\begin{proof}
Unfold the two carrier-swap relations of
\autoref{def:adjunction-unit-counit-carrier-swap-classifier}. The first gives
$p_1=q_0$, $q_1=p_0$, $a_1=a_0$, $u_1=c_0$, $c_1=u_0$, $l_1=r_0$, and
$r_1=l_0$. The second gives $p_2=q_1$, $q_2=p_1$, $a_2=a_1$, $u_2=c_1$,
$c_2=u_1$, $l_2=r_1$, and $r_2=l_1$. Substituting the first row into the
second yields
$$
p_2=p_0\land q_2=q_0\land a_2=a_0\land u_2=u_0\land c_2=c_0\land
l_2=l_0\land r_2=r_0.
$$
Substitute these coordinate identities in the carrier of $D_0$; by
\autoref{def:adjunction-unit-counit-carrier} this is the carrier of $D_2$.
The same coordinate identities turn the seven reflexive $\hsame$ witnesses on
$D_0$ into the componentwise display classifier between $D_0$ and $D_2$.
\end{proof}
\leanchecked{BEDC.Derived.AdjunctionUp.AdjunctionUnitCounitCarrier\_swap\_involution}

\begin{theorem}[Adjunction triangle carrier-swap involution]
\label{thm:adjunction-triangle-carrier-swap-involution}
Let $D_i=(p_i,q_i,a_i,u_i,c_i,l_i,r_i)$ for $i=0,1,2$ be adjunction triangle displays. If
$$
\begin{array}{l}
\mathsf{AdjunctionTriangleCarrier}(p_0,q_0,a_0,u_0,c_0,l_0,r_0)\\
{}\land \mathsf{AdjunctionUnitCounitCarrierSwap}(D_0,D_1)
\land \mathsf{AdjunctionUnitCounitCarrierSwap}(D_1,D_2),
\end{array}
$$
then the twice-swapped display is again carried and is componentwise classifier-same to the original display:
$$
\begin{array}{l}
\mathsf{AdjunctionTriangleCarrier}(p_2,q_2,a_2,u_2,c_2,l_2,r_2)\\
{}\land \mathsf{AdjunctionUnitCounitDisplaySame}(D_0,D_2).
\end{array}
$$
\end{theorem}
\begin{proof}
Unfold the first carrier-swap relation in \autoref{def:adjunction-unit-counit-carrier-swap-classifier}. It gives
$$
p_1=q_0\land q_1=p_0\land a_1=a_0\land u_1=c_0\land c_1=u_0\land l_1=r_0\land r_1=l_0.
$$
Unfold the second carrier-swap relation in the same definition. It gives
$$
p_2=q_1\land q_2=p_1\land a_2=a_1\land u_2=c_1\land c_2=u_1\land l_2=r_1\land r_2=l_1.
$$
Substituting the first row into the second row yields
$$
p_2=p_0\land q_2=q_0\land a_2=a_0\land u_2=u_0\land c_2=c_0\land l_2=l_0\land r_2=r_0.
$$
Substitute these coordinate identities in the carrier of $D_0$. By \autoref{def:adjunction-triangle-carrier}, the same two prefix-component fields and the same two continuation ledgers are exactly the carrier fields for $D_2$. The same coordinate identities also turn the seven reflexive $\hsame$ witnesses on the coordinates of $D_0$ into $\mathsf{AdjunctionUnitCounitDisplaySame}(D_0,D_2)$.
\end{proof}
